
\section{A suspension settling out through a porous bed: Wang and Ribberink's flume}

Wang and Ribberink \cite{WangRibberink1986} built an experiment whose only
sediment process is deposition. In a 30~m Delft flume 0.5~m wide, a uniform
flow 0.2155~m deep at 0.558~m/s ($u_* = 0.034$~m/s, energy slope
$0.97\times10^{-3}$) carries 100~$\mu$m sand ($w_s = 0.7$~cm/s, fed at
70.8~kg/h, below capacity so that no bedforms develop) in equilibrium over a
10~m rigid roughened bed. It then crosses onto 16~m of perforated aluminium
plate (3~mm holes, 33\,\% open) through which every grain that reaches the
bed falls into a compartmented chamber below and is not re-entrained, so the
bed boundary condition is exactly $\phi_z = -w_s c(0)$ with no pick-up. The
depth-averaged concentration decays along the test section, and the shape of
the concentration profile settles into a self-preserving form within about
2~m ($\bar u h/u_* = 3.5$~m). Eight-point siphon profiles at seven or six
stations in each of two runs are their Table~1, digitised here in full.

What makes the case worth running is that it isolates the deposition closure
and its adaptation from everything else: no entrainment, no bedload, no bed
evolution, no flow non-uniformity. \anuga{} holds the normal flow with
Dirichlet boundaries on a two-row mesh of 0.1~m cells, with the Manning $n$
that reproduces the measured $u_*$ (0.0151) and the slope normal for it, the
settling velocity set by the fraction's diameter (93~$\mu$m under
Ferguson--Church gives the measured 0.7~cm/s), the critical stress set
beyond reach so that nothing is entrained, and the inflow carrying the
measured depth mean at the first station.

The measured concentration decays exponentially from 1~m on. Writing the
decay as $\bar c \propto \exp(-\lambda w_s x/\bar u h)$, the paper's own
asymptotic theory gives $\lambda = 1.83$ to first order and 1.47 to second
order, the latter agreeing with a full two-dimensional solution by van Rijn's
SUTRENCH; the measurements give 1.44 (run~1) and 1.55 (run~2). \anuga{}'s
closures give the following, the RMS being that of
$\ln(\text{model}/\text{measured})$ over the stations from 1~m on:

\begin{center}
\footnotesize
\setlength{\tabcolsep}{4pt}
\begin{tabular}{p{4.6cm}rrrr}
\hline
closure & $\lambda$ & \multicolumn{2}{c}{ln-ratio RMS} & bias \\
 & & run 1 & run 2 & \\
\hline
measured & 1.44, 1.55 & & & \\
instantaneous exchange [D-1] (the old default) & 2.99 & 0.73 & 0.58 & $-0.52$, $-0.38$ \\
carried near-bed ratio [D-4] & 2.99 & 0.73 & 0.58 & $-0.52$, $-0.38$ \\
Galappatti lag [D-3] & 1.49 & 0.24 & 0.34 & $+0.21$, $+0.31$ \\
\textbf{two-layer, velocity split [D-5]} (the default) & \textbf{1.50} & \textbf{0.11} & \textbf{0.08} & $-0.03$, $+0.07$ \\
\hline
\end{tabular}
\end{center}

The instantaneous exchange deposits at $D = d^{*} w_s c$ with the fitted
Rouse ratio $d^{*} = 2.99$ at this Rouse number (0.50) and reference height
($0.1h$), so it decays at exactly that rate, twice as fast as the flume did:
by 16~m it has removed 93\,\% of the load where the flume removed 76\,\%.
That is the same fault the van Rijn trench shows as filling too fast, here
with no other process to hide it. The carried near-bed ratio [D-4] is
identical to it, correctly: the ratio it carries relaxes only when the flow
decelerates, and this flow is uniform. The Galappatti lag recovers the
decay rate almost exactly (its $\alpha = 1.49$ against the paper's
second-order 1.47) but decays uniformly from the transition and so misses
the sharp drop in the first metre, leaving it 20 to 30\,\% high throughout.
The two-layer suspension with the log-law velocity split gets both: the
first metre deposits fast out of the near-bed layer the equilibrium profile
arrives with, and the column then settles into the self-preserving shape and
decays at 1.50. The case runs it, and since [D-5] with the velocity split is now the
default closure the case simply takes the defaults. It requires
$\mathrm{RMS}\,\ln(\text{model}/\text{measured}) < 0.15$ from 1~m on; run~1
runs routinely and run~2 joins it under the long option.

\begin{figure}
\begin{center}
\includegraphics[width=0.95\textwidth]{concentration_run1.png}
\end{center}
\caption{Run 1: depth-averaged concentration along the test section under
each closure. The carried ratio lies under the instantaneous exchange.}
\end{figure}

\IfFileExists{concentration_run2.png}{
\begin{figure}
\begin{center}
\includegraphics[width=0.95\textwidth]{concentration_run2.png}
\end{center}
\caption{Run 2 (long run).}
\end{figure}
}{}

\endinput
